\textbf{When does deep learning pay off?} The answer depends on data and representation. At 50k
events, tuned trees are better on the richer feature sets, while the MLP is barely ahead on raw or
mass-only inputs. From 250k events upward, the MLP's ability to learn smooth, high-order interactions
dominates on every feature set, and its advantage is largest exactly where the inputs are least
processed. This reconciles the tabular-ML literature, where trees often win at typical sizes
\cite{grinsztajn2022tree,shwartzziv2022tabular}, with the HEP result that networks win on large
simulated samples \cite{baldi2014searching}. Both hold, in different regimes.

\textbf{Physics knowledge still helps the network.} Encoding the beam's rotational symmetry gives
the MLP $+0.065$ at 250k, more than the expert masses do, and at 1M it beats the raw+expert set. The
lesson is to give the network the problem's invariances rather than hand-built summaries.
Symmetry-aware architectures \cite{komiske2019energyflow} follow the same principle.

\textbf{Is the hybrid worth it?} Stacking yields the best discrimination and calibrated
probabilities at every budget, at the cost of training both models (about 27\,s at 250k vs.\ 17\,s
for the MLP alone). Its gain is largest in the mid-data regime (250k), where the two models are
comparably strong, and vanishes at 10.5M, where the MLP dominates. That is the expected behaviour of
a stacking ensemble, and the learned weights make it transparent.

\textbf{Technical validation.} The study is designed to be checked, not just run.
(i)~19 unit tests cover splits, rotation invariance, $\Delta\phi$ wrapping, metrics, DeLong
(against scikit-learn and a bootstrap), calibration, stacking and every network variant.
(ii)~The scaler, $\phi$ scale, calibrators and stacker are all fit on non-test splits.
(iii)~Seeds and subset indices are saved, and re-running a seed reproduces its numbers.
(iv)~Every claim is a paired test repeated over 3--5 seeds.

\textbf{Limitations.}
(i)~At 1M and 10.5M the MLP stopped at the 50-epoch cap (best epoch 46--50), and XGBoost at its
3000-tree cap, so large-budget results are lower bounds.
(ii)~The 10.5M run is a single seed on two feature sets. Memory limited the 47-feature sets.
(iii)~The ablations use 3 seeds at 250k only, and the dropout finding suggests a short
regularisation sweep at each budget.
(iv)~Training times above 250k were measured under memory pressure and are not compared.
(v)~The data are simulated; separability is not discovery.

\textbf{Future work.} Dropout 0.3 with a longer epoch budget at 1M+; permutation-invariant
set networks over the six physics objects; and temperature scaling inside the stacker.
